\section{Implementation Status and Baselines}

\subsection{Current Software Implementation}

The current software implementation follows the architecture described above and is maintained in the public Marine Race Arena repository~\cite{holoDroneCompetition2026}. The repository contains configuration dataclasses, race validation, arena components, simulator adapters, participant interfaces, controller loading, baseline controllers, referee state, scoring, logging, benchmark scripts, and simulator-independent tests. The implementation also includes reference tracks and command-line entry points for validating track configurations and running repeated benchmark trials.

The benchmark runner is designed to execute the same track over multiple seeds and produce per-run logs plus aggregate CSV/JSON summaries. This is important for paper-quality evaluation because single runs are not sufficient in a stochastic or disturbance-rich environment. Even deterministic controllers should be evaluated over changes in obstacle generation, current profiles, sensor noise, or initial perturbations.

\subsection{Official Baselines}

The repository currently defines two main non-cheating official baselines. The \texttt{acoustic\_baseline} is deterministic and uses beacon, sensor, and race observations. It implements an approach/transit/exit gate state machine, yaw smoothing, bearing deadbands, vertical regulation, and speed scheduling. The \texttt{rule\_gate\_baseline} combines beacon guidance with front-camera gate information when available. It uses local visual centering for yaw, sway, and heave, and reduces or stops forward surge until alignment is acceptable. Experimental or deprecated controllers, such as acoustic-vision and vision-gate variants, remain useful for debugging but should not be treated as final scientific baselines until their stability is quantified.

The key point is that these baselines are reference controllers, not the main contribution. Their role is to prove that the arena can evaluate different autonomy stacks under common rules. The final paper should eventually include completion rates and penalized times for at least acoustic-only and rule-based vision-assisted baselines on clean-gate, obstacle-gate, and current-gate tasks.

\subsection{What Is Not Yet Claimed}

This draft intentionally does not report numerical controller performance. At the time of writing, the current-gate baseline with marine currents is still under development, and no full benchmark table has been finalized. Table~\ref{tab:implementation_status} records the current state without fabricating results.

\begin{table}[t]
\centering
\scriptsize
\renewcommand{\arraystretch}{1.08}
\caption{Current implementation and evidence status.}
\label{tab:implementation_status}
\begin{tabular}{@{}p{0.25\linewidth}p{0.34\linewidth}p{0.23\linewidth}@{}}
\toprule
Component & Current status & Quantitative evidence \\
\midrule
Configurable race tracks & Declarative track files & Structural validation \\
Gate referee & Center-point aperture crossing & Unit validation \\
Scoring and logs & Event logs and summaries & Needs benchmark runs \\
Acoustic baseline & Official non-cheating controller & Needs task table \\
Rule gate baseline & Implemented; robust transit tuning & Needs task table \\
Current-gate task & Defined in configuration & Tuning ongoing \\
Multi-ROV task & Future-ready task/interface & No final experiment \\
Real-world validation & Not available in this draft & Not claimed \\
\bottomrule
\end{tabular}
\end{table}

\subsection{Recommended Final Experimental Matrix}

The final experimental section should contain, at minimum, three controllers and three benchmark tasks. The controllers should be: an acoustic-only baseline, the rule-based gate baseline, and an oracle/debug controller used only as an upper-bound track sanity check. The tasks should be clean-gate, obstacle-gate, and current-gate. For scientific reporting, each controller-task pair should be executed over multiple seeds, with fixed random seeds recorded in metadata. The oracle should be excluded from official ranking tables or clearly marked as debug-only.

A complete result table should include completion rate, mean completed gates, mean official time for completed runs, mean penalized time, collisions, out-of-bounds events, stuck events, and DNF reasons. Only after this table is produced should the abstract and conclusion contain performance claims.
